% =============================================================
% Annexure I: Software and Environment Details
% =============================================================
\section{Software and Environment Details}
\label{annexure:software}


% -------------------------------------------------------------
\subsection{Python Distribution and Version}
\label{annexure:python}

Python~3.11 via the Anaconda distribution. Jupyter Notebook was used for interactive development.

% -------------------------------------------------------------
\subsection{Core Package Versions}
\label{annexure:packages}

\begin{table}[htbp]
\centering
\caption{Python package versions used in the dissertation.}
\label{tab:packages}
\begin{tabular}{lll}
\hline
\textbf{Package} & \textbf{Version} & \textbf{Purpose} \\
\hline
PyTorch           & 2.x              & DNN architecture, training and inference \\
SHAP              & 0.43+            & Shapley value computation (GradientExplainer) \\
statsmodels       & 0.14+            & GLM (Gamma, log-link) fitting \\
NumPy             & 1.24+            & Numerical computation \\
pandas            & 2.x              & Data manipulation and I/O \\
SciPy             & 1.11+            & Statistical tests, distributions \\
scikit-learn      & 1.3+             & GP surrogate, train/test splitting, metrics \\
matplotlib        & 3.8+             & Static figure generation \\
seaborn           & 0.13+            & Statistical visualisation \\
scikit-optimize   & 0.10+            & Bayesian Optimisation (\texttt{gp\_minimize}) \\
Plotly            & 5.x              & Interactive 3D surface plots \\
\hline
\end{tabular}

\medskip
\emph{Note:} Exact patch versions may vary. The versions above represent the minimum versions with which the code was tested. Update these entries with the output of \texttt{pip freeze} or \texttt{conda list} before final submission.
\end{table}

% -------------------------------------------------------------
\subsection{Hardware Platform}
\label{annexure:hardware}

\begin{itemize}
    \item Processor: Apple Silicon (M-series)
    \item Primary compute device: MPS (Metal Performance Shaders) via PyTorch
    \item SHAP compute device: CPU (faster than MPS for GradientExplainer)
    \item Fallback: CPU (used when MPS is unavailable or for SHAP computations)
\end{itemize}

\begin{lstlisting}[caption={Device selection logic for PyTorch.}]
import torch

if torch.backends.mps.is_available():
    DEVICE = torch.device('mps')
elif torch.cuda.is_available():
    DEVICE = torch.device('cuda')
else:
    DEVICE = torch.device('cpu')
\end{lstlisting}

% -------------------------------------------------------------
\subsection{Reproducibility}
\label{annexure:reproducibility}

Fixed random seed of 42:
\begin{itemize}
    \item PyTorch weight initialisation (\texttt{torch.manual\_seed(42)})
    \item NumPy random operations (\texttt{np.random.seed(42)})
    \item NumPy Generator objects (\texttt{np.random.default\_rng(42)})
    \item scikit-learn train/test split (\texttt{random\_state=42})
    \item GP surrogate fitting (\texttt{random\_state=42})
\end{itemize}

% -------------------------------------------------------------
\subsection{Data Files}
\label{annexure:data_files}

\begin{itemize}
    \item \texttt{insurance\_dataset.csv}: 1\,000\,000 rows, 12 columns (age, gender, bmi, children, smoker, region, medical\_history, family\_medical\_history, exercise\_frequency, occupation, coverage\_level, charges). Synthetic benchmark dataset sourced from Kaggle.

    \item \texttt{chapter3\_final\_model.pth}: PyTorch checkpoint containing the trained FunnelDNN model state, architecture specification, hyperparameters, standardisation parameters, feature names and evaluation metrics.

    \item \texttt{model\_ready\_charges\_train\_20260517\_115748.csv}: reference-category encoded training partition (60\% of the dataset, 22 features + charges).
\end{itemize}
